\section{Phép tính luỹ thừa}

\subsection{Đề 1}
\Opensolutionfile{ans}[ans/ans-0-B1]
\caulc
\begin{ex}%%[1D6N1-3]
	Nếu $\left(a-2\right)^{\tfrac{1}{4}}<\left(a-2\right)^{\tfrac{1}{3}}$ thì khẳng định nào sau đây là đúng?
	\choice
	{$2 < a < 3$}
	{$a > 2$}
	{$a < 3$}
	{\True $a > 3$}
	\loigiai{
		Ta có: $\dfrac{1}{4}<\dfrac{1}{3}$ và $\left(a-2\right)^{\tfrac{1}{4}}<\left(a-2\right)^{\tfrac{1}{3}}$ nên $a-2 > 1\Leftrightarrow a > 3$.}
\end{ex}

\begin{ex}%[1D6N1-1]
	Với $a$ là số thực dương tùy ý, $\sqrt[14]{a}$ bằng
	\choice
	{\True $a^{\tfrac{1}{14}}$}
	{$a^{\sqrt{14}}$}
	{$a^{14}$}
	{$\sqrt{a^{14}}$}
	\loigiai{
		Ta có $\sqrt[14]{a} = a^{\tfrac{1}{14}}$.
	}
\end{ex}

\begin{ex}%[1D6N1-1]
	Với $a$ là số thực dương tùy ý, tích $a^2\cdot a^{\tfrac{1}{3}}$ bằng
	\choice
	{\True $a^{\tfrac{7}{3}}$}
	{$a^{\tfrac{2}{3}}$}
	{$a^{\tfrac{5}{3}}$}
	{$a^{\tfrac{4}{3}}$}
	\loigiai{
		Ta có $a^2\cdot a^{\tfrac{1}{3}}= a^{\tfrac{7}{3}}$.
	}
\end{ex}

\begin{ex}%[1D6N1-1]
	Cho các số dương $a\ne 1$ và các số thực $\alpha $, $\beta $. Đẳng thức nào sau đây là sai?
	\choice
	{$a^\alpha\cdot a^\beta=a^{\alpha+\beta}$}
	{\True $a^\alpha\cdot a^\beta=a^{\alpha\beta}$}
	{$\dfrac{a^\alpha}{a^\beta}=a^{\alpha-\beta}$}
	{$\left(a^\alpha\right)^\beta=a^{\alpha\beta}$}
	\loigiai{
	}
\end{ex}

\begin{ex}%[1D6H1-2]
	Cho $x > 0$. Viết biểu thức $\sqrt[5]{x^7}:\sqrt{x}$ dưới dạng lũy thừa với số mũ hữu tỉ.
	\choice
	{$x^{\tfrac{9}{7}}$}
	{$x^{\tfrac{9}{5}}$}
	{$x^{\tfrac{7}{10}}$}
	{\True $x^{\tfrac{9}{10}}$}
	\loigiai{
		Ta có  $\sqrt[5]{x^7}:\sqrt{x}=x^{\tfrac{7}{5}}:x^{\tfrac{1}{2}}=x^{\tfrac{9}{10}}$.
	}
\end{ex}

\begin{ex}%[1D6H1-2]
	Viết biểu thức $P=\dfrac{a^2a^{\tfrac{5}{2}}\sqrt[3]{a^4}}{\sqrt[6]{a^7}}$, $\left(a > 0\right)$ dưới dạng lũy thừa với số mũ hữu tỉ.
	\choice
	{$P=a^5$}
	{\True $P=a^{\tfrac{14}{3}}$}
	{$P=a^4$}
	{$P=a^{\tfrac{13}{3}}$}
	\loigiai{
		Ta có $a^2a^{\tfrac{5}{2}}\sqrt[3]{a^4}=a^2a^{\tfrac{5}{2}}a^{\tfrac{4}{3}} = a^{\tfrac{35}{6}}$ và $\sqrt[6]{a^7}=a^{\tfrac{7}{6}}$.\\
		Do đó $P = a^{\tfrac{14}{3}}$.
	}
\end{ex}

\begin{ex}%[1D6H1-4]
	So sánh ba số: $\left(0,2\right)^{0,3}$, $\left(0,7\right)^{3,2}$ và $\left(\sqrt 3\right) ^{0,2}$ ta được
	\choice
	{\True $\left(0,7\right)^{3,2}<\left(0,2\right)^{0,3}<\left(\sqrt 3\right)^{0,2}$}
	{$\left(0,2\right)^{0,3}<\left(0,7\right)^{3,2}<\left(\sqrt 3\right)^{0,2}$}
	{$\left(\sqrt 3\right) ^{0,2}<\left(0,2\right)^{0,3}<\left(0,7\right)^{3,2}$}
	{$\left(0,2\right)^{0,3}<\left(\sqrt 3\right)^{0,2}<\left(0,7\right)^{3,2}$}
	\loigiai{
		Ta có $\left(0,2\right)^{0,3}=\left(0,2\right)^{\tfrac{3}{10}}=\left[\left(0,2\right)^3\right]^{\tfrac{1}{10}}=\left(0,008\right)^{\tfrac{1}{10}}$.\\
		Mặt khác
		$\left(0,7\right)^{3,2}=\left(0,7\right)^{\tfrac{32}{10}}=\left[\left(0,7\right)^{32}\right]^{\tfrac{1}{10}}$ 
		và 
		$\left(\sqrt 3\right) ^{0,2}=(3)^{\tfrac{1}{2}.\tfrac{2}{10}}=3^{\tfrac{1}{10}}$.\\
		Do $\left(0,7\right)^{32}< 0,008 < 3$ nên $\left(0,7\right)^{3,2}<\left(0,2\right)^{0,3}<\left(\sqrt 3\right) ^{0,2}$.}
\end{ex}

\begin{ex}%[1D6H1-2]
	Biểu thức $K=\sqrt[3]{\dfrac{2}{3}\sqrt[3]{\dfrac{2}{3}\sqrt{\dfrac{2}{3}}}}$ viết dưới dạng luỹ thừa với số mũ hữu tỉ là
	\choice
	{$\left(\dfrac{2}{3}\right)^{\tfrac{5}{18}}$}
	{\True $\left(\dfrac{2}{3}\right)^{\tfrac{1}{2}}$}
	{$\left(\dfrac{2}{3}\right)^{\tfrac{1}{8}}$}
	{$\left(\dfrac{2}{3}\right)^{\tfrac{1}{6}}$}
	\loigiai{
		Coi $X=\dfrac{2}{3}$. Theo nguyên tắc \lq\lq Chia cộng\rq\rq \,ta có\\
		$K=\sqrt[3]{X\sqrt[3]{X\sqrt X}}=\sqrt[3]{X\sqrt[3]{X.X^{\tfrac{1}{2}}}}=\sqrt[3]{X\sqrt[3]{X^{\tfrac{3}{2}}}}=\sqrt[3]{X.X^{\tfrac{1}{2}}}=\sqrt[3]{X^{\tfrac{3}{2}}}=X^{\tfrac{1}{2}}$.
	}
\end{ex}

\begin{ex}%[1D6H1-2]
	Cho biểu thức $P=\sqrt x \cdot \sqrt[3]{x}\cdot\sqrt[6]{x^4}$, $\left(x > 0\right)$. Mệnh đề nào dưới đây đúng?
	\choice
	{$P=x^{\tfrac{2}{3}}$}
	{$P=x^{\tfrac{5}{2}}$}
	{\True $P=x^{\tfrac{3}{2}}$}
	{$P=x^{\tfrac{2}{5}}$}
	\loigiai{
		Ta có $P=\sqrt x \cdot \sqrt[3]{x}\cdot\sqrt[6]{x^4} = x^{\tfrac{1}{2}}\cdot x^{\tfrac{1}{3}}\cdot \tfrac{2}{3}=x^{\tfrac{3}{2}}$.
	}
\end{ex}

\begin{ex}%[1D6H1-2]
	Cho $x > 0$. Viết biểu thức $\sqrt[5]{x^7}:\sqrt x $ dưới dạng lũy thừa với số mũ hữu tỉ.
	\choice
	{$x^{\tfrac{9}{7}}$}
	{$x^{\tfrac{9}{5}}$}
	{$x^{\tfrac{7}{10}}$}
	{\True $x^{\tfrac{9}{10}}$}
	\loigiai{
		Ta có $\sqrt[5]{x^7}:\sqrt x = x^{\tfrac{7}{5}}\cdot x^{\tfrac{1}{2}}=x^{\tfrac{9}{10}}$.
	}
\end{ex}

\begin{ex}%[1D6H1-1]
	Cho $a,\,b$ là các số thực dương và thỏa mãn $a=5^x$, $b=3^x$. Giá trị của biểu thức $P=25^x+15^x+27^x$ bằng
	\choice
	{$P=a^3+ab+b^2$}
	{$P=a^2b^2+ab+b^2$}
	{\True $P=a^2+ab+b^3$}
	{$P=a{b^3}+ab+b^2$}
	\loigiai{
		Ta có $P=25^x+15^x+27^x=\left(5^2\right)^x+\left(5.3\right)^x+\left(3^3\right)^x=\left(5^x\right)^2+5^x{3^x}+\left(3^x\right)^3$.\\
		Với $a=5^x$, $b=3^x$ thì $P=a^2+ab+b^3$.
	}
\end{ex}

\begin{ex}%[1D6H1-1]
	Rút gọn biểu thức $A=\dfrac{\sqrt[3]{a^5}\cdot a^{\tfrac{7}{3}}}{a^4\cdot\sqrt[7]{a^{-2}}}$ với $a > 0$ ta được kết quả $A=a^{\tfrac{m}{n}}$, trong đó $m, n\in\mathbb{N}^*$ và $\dfrac{m}{n}$ là phân số tối giản. Khẳng định nào sau đây đúng?
	\choice
	{$3m^2-2n=2$}
	{$m^2+n^2=43$}
	{\True $2m^2+n=15$}
	{$m^2+n^2=25$}
	\loigiai{
		Ta có $A=\dfrac{\sqrt[3]{a^5}\cdot a^{\tfrac{7}{3}}}{a^4\cdot\sqrt[7]{a^{-2}}}=\dfrac{a^{\tfrac{5}{3}}\cdot a^{\tfrac{7}{3}}}{a^4\cdot a^{-\tfrac{2}{7}}}=\dfrac{a^{\tfrac{5}{3}+\tfrac{7}{3}}}{a^{4-\tfrac{2}{7}}}=\dfrac{a^4}{a^{\tfrac{26}{7}}}=a^{4-\tfrac{26}{7}}=a^{\tfrac{2}{7}}$.\\
		Suy ra $\left\{\begin{array}{l}
			m=2\\
			n=7
		\end{array}\right.\Rightarrow 2m^2+n=2\cdot2^2+7=15$.}
\end{ex}
\Closesolutionfile{ans}
\indapan{6}{ans/ans-0-B1}
\cauds
\Opensolutionfile{ans}[ans/ans-0-B1-DS]

\begin{ex}
	Cho biểu thức $ A=\left(\dfrac{a^{\sqrt 5}}{b^{\sqrt 5-2}}\right)^{\sqrt 5+2}\cdot\dfrac{a^{-2-\sqrt 5}}{b^{-1}}$ với $ a,b > 0$. Vậy
	\choiceTF
	{Sau khi rút gọn, thì biểu thức $ A$ chỉ chứa biến $ b$}
	{Với $ a=2,b=1+5\sqrt 2 $ thì $ A=\dfrac{113}{3}$}
	{\True Khi $ A=a^m\cdot b^n$ thì $ m+n=3+\sqrt 5 $}
	{Khi $ A=a^m\cdot b^n$ thì $ m-n=2+\sqrt 5 $}
	\loigiai{
		Ta có $ A=\dfrac{\left(a^{\sqrt 5}\right)^{\sqrt 5+2}}{\left(b^{\sqrt 5-2}\right)^{\sqrt 5+2}}\cdot\dfrac{b}{a^{2+\sqrt 5}}=\dfrac{a^{5+2\sqrt 5}\cdot b}{b\cdot{a^{2+\sqrt 5}}}=a^{3+\sqrt 5}$.
		\begin{itemchoice}
			\itemch Sai. Vì $A =a^{3+\sqrt 5}$.
			\itemch Sai. Vì $ a=2,b=1+5\sqrt 2$ thì $ A=2^{3+\sqrt 5}$.
			\itemch Đúng. Vì $A=a^m\cdot b^n=a^{3+\sqrt 5}\Leftrightarrow \heva{&m=3+\sqrt 5\\&n=0.}$\\
			Vậy $m+n=3+\sqrt 5$.
			\itemch Sai. Vì $m-n=3+\sqrt 5$.
		\end{itemchoice}
	}
\end{ex}

\begin{ex}
	Cho biểu thức $ B=\dfrac{a^{\tfrac{1}{3}}\left(a^{\tfrac{1}{2}}-a^{\tfrac{5}{2}}\right)}{a^{\tfrac{1}{4}}\left(a^{\tfrac{7}{12}}-a^{\tfrac{19}{12}}\right)}$ với $ a > 0$.
	\choiceTF
	{\True Sau khi rút gọn, thì biểu thức $ B=k+a^m,\left(k,m > 0\right)$}
	{Với $ a=4+\sqrt 5 $ thì $ B=6+3\sqrt 5 $}
	{\True $\dfrac{B}{1-a^2}=\dfrac{1}{1-a}$}
	{Với $a_1, a_2$ là $2$ nghiệm của $B^2=4$ thì $\dfrac{1}{x_1}+\dfrac{1}{x_2}=23$}
	\loigiai{
		Ta có $ B=\dfrac{a^{\tfrac{1}{3}}\left(a^{\tfrac{1}{2}}-a^{\tfrac{5}{2}}\right)}{a^{\tfrac{1}{4}}\left(a^{\tfrac{7}{12}}-a^{\tfrac{19}{12}}\right)}=\dfrac{a^{\tfrac{1}{3}}\cdot{a^{\tfrac{1}{2}}}\left(1-a^2\right)}{a^{\tfrac{1}{4}}\cdot{a^{\tfrac{7}{12}}}(1-a)}=\dfrac{a^{\tfrac{5}{6}}(1+a)}{a^{\tfrac{5}{6}}}=1+a$.
		\begin{itemchoice}
			\itemch Đúng. Vì $B = 1+a$.
			\itemch Sai. Với $ a=4+\sqrt 5 $ thì $ B=5+\sqrt 5 $.
			\itemch Đúng. Vì  $\dfrac{B}{1-a^2}=\dfrac{1+a}{\left(1-a\right)\left(1+a\right)}=\dfrac{1}{1-a}$.
			\itemch Sai. Vì $\left(a+1\right)^2=4  \Leftrightarrow a^2+2a-3=0$.\\
			Vậy $\dfrac{1}{a_1}+\dfrac{1}{a_2}=\dfrac{2}{3}$.
		\end{itemchoice}
	}
\end{ex}

\begin{ex}
	Cho biểu thức $ P=\dfrac{a^{\tfrac{1}{3}}\sqrt b+b^{\tfrac{1}{3}}\sqrt a}{\sqrt[6]{a}+\sqrt[6]{b}}-\sqrt[3]{ab}$ với $ a,b > 0$.
	\choiceTF
	{$ P=a+2b$}
	{Với $ a=\sqrt 5,b=\sqrt 3 $ thì $ P=\sqrt 5+2\sqrt 3 $}
	{\True $ P=k$ ($ k$ là hằng số)}
	{\True Với $ a=\sqrt{22}$, $b=4$ thì $ P=0$}
	\loigiai{
		Ta có
		\begin{eqnarray*}
			P&=&\dfrac{a^{\tfrac{1}{3}}\sqrt b+b^{\tfrac{1}{3}}\sqrt a}{\sqrt[6]{a}+\sqrt[6]{b}}-\sqrt[3]{ab}\\
			&=&\dfrac{a^{\tfrac{1}{3}}{b^{\tfrac{1}{2}}}+b^{\tfrac{1}{3}}{a^{\tfrac{1}{2}}}}{a^{\tfrac{1}{6}}+b^{\tfrac{1}{6}}}-(ab)^{\tfrac{1}{3}}\\
			&=&\dfrac{a^{\tfrac{1}{3}}{b^{\tfrac{1}{3}}}\left(b^{\tfrac{1}{6}}+a^{\tfrac{1}{6}}\right)}{a^{\tfrac{1}{6}}+b^{\tfrac{1}{6}}}-(ab)^{\tfrac{1}{3}}\\
			&=&a^{\tfrac{1}{3}}{b^{\tfrac{1}{3}}}-(ab)^{\tfrac{1}{3}}=0.
		\end{eqnarray*}
		\begin{itemchoice}
			\itemch Sai. Vì $P = 0$.
			\itemch Sai. Vì $P = 0$.
			\itemch Đúng. Vì $P = 0$.
			\itemch Đúng. Vì $P = 0$.
		\end{itemchoice}
	}
\end{ex}

\begin{ex}
	Cho các biểu thức $ A=\sqrt{2\cdot\sqrt[3]{2\cdot\sqrt[4]{2}}},\,B=\sqrt[24]{2^5}\cdot\dfrac{1}{\sqrt{2^{-1}}}$. Vậy
	\choiceTF
	{\True $ A=2^{\tfrac{a}{b}}$($\dfrac{a}{b}$ là phân số tối giản), khi đó $ a+b=41$}
	{$ B=2^{\tfrac{a}{b}}$($\dfrac{a}{b}$ là phân số tối giản), khi đó $ a+b=31$}
	{$ A-B\sqrt 5=\sqrt 5 $}
	{\True $ A\cdot B=2^{\tfrac{m}{n}}$($\dfrac{m}{n}$ là phân số tối giản), khi đó $ m+n=29$}
	\loigiai{
		Ta có
		\begin{eqnarray*}
			A&=&\sqrt{2\cdot\sqrt[3]{2\cdot\sqrt[4]{2}}}\\
			&=&\sqrt{2\sqrt[3]{2\cdot{2^{\tfrac{1}{4}}}}}\\
			&=&\sqrt{2\cdot{2^{\tfrac{5}{12}}}}\\
			&=&2^{\tfrac{17}{24}}.
		\end{eqnarray*}
		và $B = \sqrt[24]{2^5}\cdot\dfrac{1}{\sqrt{2^{-1}}}=2^{\tfrac{5}{24}}\cdot{2^{\tfrac{1}{2}}}=2^{\tfrac{17}{24}}$.
		\begin{itemchoice}
			\itemch Đúng. Vì $A = 2^{\tfrac{17}{24}}$ nên $a+b=41$.
			\itemch Sai. Vì $B = 2^{\tfrac{17}{24}}$ nên $a+b=41$.
			\itemch Sai. Vì $A-B\sqrt 5=\sqrt 5 = 2^{\tfrac{17}{24}}(1-\sqrt 5)$.
			\itemch Đúng. $ A\cdot B= 2^{\tfrac{34}{24}}$. khi đó $ m+n=29$.
		\end{itemchoice}
	}
\end{ex}
\Closesolutionfile{ans}
\indapan{3}{ans/ans-0-B1-DS}

\Opensolutionfile{ans}[ans/ans-0-B1-KQ]
\caukq
\begin{ex}%[1D6H1-1]
	Tính giá trị của biểu thức $ P=\dfrac{2^3\cdot{2^{-1}}+5^{-3}\cdot{5^4}}{10^{-3}:{10^{-2}}-0{,}1^0}$.
	\shortans{$-10$}
	\loigiai{
		Ta có $ P=\dfrac{2^3\cdot{2^{-1}}+5^{-3}\cdot{5^4}}{10^{-3}:{10^{-2}}-0,1^0}=\dfrac{2^2+5}{10^{-1}-1}=-10$.}
\end{ex}

\begin{ex}%[1D6H1-1]
	Cho biểu thức sau $B=\dfrac{a^{\tfrac{4}{3}}\left(a^{-\tfrac{1}{3}}+a^{\tfrac{2}{3}}\right)}{a^{\tfrac{1}{4}}\left(a^{\tfrac{3}{4}}+a^{-\tfrac{1}{4}}\right)}$ với $a>0$. Tính $B$ khi $a = 100$.
	\shortans{$100$}
	\loigiai{
		Ta có $B=\dfrac{a^{\tfrac{4}{3}}\left(a^{-\tfrac{1}{3}}+a^{\tfrac{2}{3}}\right)}{a^{\tfrac{1}{4}}\left(a^{\tfrac{3}{4}}+a^{-\tfrac{1}{4}}\right)}=\dfrac{a+a^2}{a+1}=\dfrac{a(a+1)}{a+1}=a$.\\
		Vậy $B = 100 = 100$.
	}
\end{ex}

\begin{ex}%[1D6H1-1]
	Biết $4^x+4^{-x}=23$, tính giá trị biểu thức $ P=2^x+2^{-x}$.
	\shortans{$5$} 
	\loigiai{
		Đặt $ P=2^x+2^{-x}\Rightarrow P > 0$.\\
		Ta có
		\begin{eqnarray*}
			P^2&=&\left(2^x+2^{-x}\right)^2\\
			&=&4^x+4^{-x}+2\cdot{2^x}\cdot{2^{-x}}\\
			&=& 23+2=25.
		\end{eqnarray*}
		Do đó $ P=5$.
	}
\end{ex}

\begin{ex}%[1D6H3-5]
	Giả sử số tiền gốc là $ A$, lãi suất là $ r\% /$ kì hạn gửi (có thể là tháng, quý hay năm) thì tồng số tiền nhận được cả gốc và lãi sau $ n$ kì hạn gửi là $ A{(1+r)^n}$. Bà Hạnh gửi 50 triệu vào tài khoản định kỳ tính lãi kép với lãi suất là $ 8\% /$ năm. Tính số tiền lãi thu được sau 10 năm (làm tròn đến hàng phần mười).
	\shortans{$57{,}9$}
	\loigiai{
		Áp dụng công thức tính lãi kép, sau 10 năm số tiền cả gốc và lãi bà Hạnh thu về là $ A{(1+r)^n}=50(1+0,08)^{10}\approx 107,9$ triệu đồng.\\
		Suy ra số tiền lãi bà Hạnh thu về sau 10 năm là $ 107{,}9-50=57{,}9$ triệu đồng.}
\end{ex}

\begin{ex}%[1D6H3-5]
	Số lượng của loại vi khuẩn A trong một phòng thí nghiệm được tính theo công thức $ s(t)=s(0)\cdot{2^t}$, trong đó $ s(0)$ là số lượng vi khuẩn $ A$ lúc ban đầu, $ s(t)$ là số lượng vi khuẩn $ A$ có sau $ t$ phút. Biết sau 3 phút thì số lượng vi khuẩn $ A$ là 625 con. Hỏi sau bao nhiêu phút, kể từ lúc ban đầu, số lượng vi khuẩn A là 10 triệu con?\\
	\shortans{$7$}
	\loigiai{
		Ta có $ s(3)=s(0)\cdot{2^3}\Rightarrow s(0)=\dfrac{s(3)}{8}=78,125$ nghìn con.\\
		Do đó $ s(t)=10$ triệu con $=10000$ nghìn con khi
		\begin{eqnarray*}
			& &10000=s(0)\cdot{2^t} \\
			&\Leftrightarrow& {2^t}=\dfrac{10000}{78,125}=128=2^7\\
			&\Leftrightarrow & {2^t}=2^7\\
			&\Leftrightarrow & t=7.
		\end{eqnarray*}
	}
\end{ex}

\begin{ex}%[1D6V1-1]
	Biết $9^\alpha=\dfrac{1}{2}$. Tính $ B=\left(3^\alpha+3^{-\alpha}\right)^2-\left(81^\alpha+81^{-\alpha}\right)$.\\
	\shortans{$8{,}25$}
	\loigiai{
		Ta có
		\begin{eqnarray*}
			B&=& \left(3^\alpha\right)^2+2\cdot{3^\alpha}\cdot{3^{-\alpha}}+\left(3^{-\alpha}\right)^2-\left(9^2\right)^\alpha+\left(9^2\right)^{-\alpha}\\
			&=& 3^{2\alpha}+2\cdot{3^{\alpha+(-\alpha)}}+\left(3^2\right)^{-\alpha}-\left(9^\alpha\right)^2+\left(9^{-\alpha}\right)^2\\
			&= & 9^\alpha+2\cdot{3^0}+9^{-\alpha}-\left(9^\alpha\right)^2+\left(9^\alpha\right)^{-2}\\
			&= & \dfrac{1}{2}+2+\left(\dfrac{1}{2}\right)^{-1}-\left(\dfrac{1}{2}\right)^2+\left(\dfrac{1}{2}\right)^{-2}\\
			&=&8{,}25.
		\end{eqnarray*}
	}
\end{ex}

\Closesolutionfile{ans}
\indapan{6}{ans/ans-0-B1-KQ}
\subsection{Đề 2}
\Opensolutionfile{ans}[ans/ans-0-B1]
\caulc
\begin{ex}%[1D6N1-3]
	Cho $a$, $b>0$; $\alpha$, $\beta \in \mathbb{R}$. Mệnh đề nào sau đây \textbf{sai}?
	\choice
	{$\left(a^{\alpha}\right)^{\tfrac{1}{\beta}}=a^{\tfrac{\alpha }{\beta }}$, $\beta \neq 0$}
	{$(a\cdot b)^{\alpha }=a^{\alpha}\cdot b^{\alpha}$}
	{$\dfrac{a^{\alpha}}{a^{\beta}}=a^{\alpha-\beta}$}
	{\True $a^{\alpha}\cdot b^{\beta}=(ab)^{\alpha +\beta }$}
	\loigiai{
		Với $a$, $b>0$ và $\alpha$, $\beta \in \mathbb{R}$ ta có khẳng định $a^{\alpha}\cdot b^{\beta}=(ab)^{\alpha +\beta }$ sai, các khẳng định còn lại đúng.
	}
	
\end{ex}


\begin{ex}%[1D6H1-1]
	Cho $a$ là số thực dương. Giá trị rút gọn của biểu thức $P=a^{\tfrac{4}{3}} \sqrt{a}$ bằng
	\choice
	{\True$a^{\tfrac{11}{6}}$}
	{$a^{\tfrac{10}{3}}$}
	{$a^{\tfrac{7}{3}}$}
	{$a^{\tfrac{5}{6}}$}
	\loigiai{
		Ta có $P=a^{\tfrac{4}{3}} \sqrt{a}=a^{\tfrac{4}{3}} \cdot a^{\tfrac{1}{2}}=a^{\tfrac{4}{3}+\tfrac{1}{2}}=a^{\tfrac{11}{6}}$.}
	
\end{ex}


\begin{ex}%[1D6H1-2]
	Cho $a$ là số thực dương khác $1$. Giá trị của biểu thức $P=a^{\tfrac{2}{3}} \sqrt{a}$ bằng
	\choice
	{$a^{\tfrac{2}{3}}$}
	{$a^{\tfrac{5}{6}}$}
	{$a^3$}
	{\True $a^{\tfrac{7}{6}}$}
	\loigiai
	{
		Ta có $P=a^{\tfrac{2}{3}} \sqrt{a}=a^{\tfrac{2}{3}} \cdot a^{\tfrac{1}{2}}=a^{\tfrac{7}{6}}$.
	}
\end{ex}


\begin{ex}%[1D6H1-1]
	Cho $P=\dfrac{a \sqrt{a} \sqrt[3]{a^2}}{(\sqrt[4]{a})^3}$, ký hiệu $x=\sqrt[12]{a}$. Hãy biễu diễn $P$ theo $x$.
	\choice
	{$P=x^{12}$}
	{$P=x^{10}$}
	{\True $P=x^{17}$}
	{$P=x^{\tfrac{17}{12}}$}
	\loigiai{
		Ta có $P=\dfrac{a\sqrt{a} \sqrt[3]{a^2}}{(\sqrt[4]{a})^3}=\dfrac{a^{\tfrac{13}{6}}}{a^{\tfrac{3}{4}}}= a^{\tfrac{17}{12}}$.\\
		Mặt khác $x=\sqrt[12]{a}\Leftrightarrow a = x^{12}$.\\
		Vậy $P = x^{17}$.
	}
\end{ex}

\begin{ex}%[1D6N1-1]
	Với $a$ là số thực dương tùy ý, $(\sqrt[3]{a})^{10}: a^2=a^{\tfrac{p}{q}}$ với $p, q \in \mathbb{Z}$ và $\dfrac{p}{q}$ là phân số tối giản. Giá trị của $p+q$ bằng
	\choice
	{$23$}
	{\True $7$}
	{$8$}
	{$19$}
	\loigiai{
		Ta có $(\sqrt[3]{a})^{10}: a^2=a^{\tfrac{4}{3}}$.\\
		Vậy $p+q = 7$.
	}
\end{ex}
\begin{ex}%[1D6N1-4]
	Cho số thực dương $a$ thỏa mãn $\left( \dfrac{1}{a}\right)^{24}>\left( \dfrac{1}{a}\right)^{23} $. Khẳng định nào sau đây đúng?
	\choice
	{\True $0<a<1$}
	{$a<\dfrac{1}{2}$}
	{$\dfrac{1}{2}<a<1$}
	{$a>1$}
	\loigiai{
		Với $a>0$ và $\left( \dfrac{1}{a}\right)^{24}>\left( \dfrac{1}{a}\right)^{23} $, ta suy ra $\dfrac{1}{a}>1$ hay $a<1$.\\
		Vậy $0<a<1$.
	}
	
\end{ex}


\begin{ex}%[1D6H1-2]
	Cho biểu thức $P=x^{-\tfrac{3}{4}}\cdot\sqrt{\sqrt{x^5}}$ với $x>0$. Khẳng định nào sau đây là đúng?
	\choice
	{$x^{-2}$}
	{$x^2$}
	{$x^{-\tfrac{1}{2}}$}
	{\True $x^{\tfrac{1}{2}}$}
	\loigiai{
		Ta có $P=x^{-\tfrac{3}{4}}\cdot x^{\tfrac{5}{4}}=x^{-\tfrac{3}{4}+\tfrac{5}{4}}=x^{\tfrac{1}{2}}$.}
	
\end{ex}


\begin{ex}%[1D6N1-1]
	Cho các số thực $a$, $b$, $n$, $m$ ($a,b>0$). Khẳng định nào sau đây là đúng?
	\choice
	{$\left( a^m\right)^n=a^{m+n} $}
	{$(a+b)^m=a^m+b^m$}
	{$\dfrac{a^m}{a^n}=\sqrt[n]{a^m}$}
	{\True $a^m\cdot a^n=a^{m+n}$}
	\loigiai{
		Khẳng định đúng là $a^m\cdot a^n=a^{m+n}$.
	}
	
\end{ex}


\begin{ex}%[1D6N1-2]
	Với $a$ là số thực dương tùy ý, $\sqrt{a^3}$ bằng
	\choice
	{\True $a^{\tfrac{3}{2}}$}
	{$a^6$}
	{$a^{\tfrac{1}{6}}$}
	{$a^{\tfrac{2}{3}}$}
	\loigiai{
		Với $a>0$ ta có $\sqrt{a^3}=a^{\tfrac{3}{2}}$.}
\end{ex}
\begin{ex}%[1D6N1-1]
	Thu gọn biểu thức $A=\sqrt[4]{a^3} \cdot a$ với $a$ là số thực dương ta được
	\choice
	{$A=a^{\tfrac{5}{2}}$}
	{\True $A=a^{\tfrac{7}{4}}$}
	{$A=a^{\tfrac{3}{4}}$}
	{$A=a^{\tfrac{1}{4}}$}
	\loigiai{
		Ta có $A=\sqrt[4]{a^3} \cdot a=a^{\tfrac{3}{4}}\cdot a=a^{\tfrac{7}{4}}$.
	}
	
\end{ex}


\begin{ex}%[1D6H1-2]
	Cho biểu thức $P=x^{\tfrac{1}{2}} \cdot x^{\tfrac{1}{3}} \cdot \sqrt[6]{x}$ với $x > 0$. Mệnh đề nào dưới đây đúng?
	\choice
	{$P=x^{\tfrac{5}{6}}$}
	{$P=x^{\tfrac{11}{6}}$}
	{$P=x^{\tfrac{7}{6}}$}
	{\True $P=x$}
	\loigiai{
		Ta có $P=x^{\tfrac{1}{2}} \cdot x^{\tfrac{1}{3}}\cdot x^{\tfrac{1}{6}}=x^{\tfrac{1}{2}+\tfrac{1}{3}+\tfrac{1}{6}}=x$.
	}
\end{ex}



\begin{ex}%[1D6N1-2]
	Với $a$ là số thực dương tùy ý, biểu thức $a^{\tfrac{5}{3}}\cdot a^{\tfrac{1}{3}}$ là
	\choice
	{$a^{5}$}
	{$a^{\tfrac{5}{9}}$}
	{$a^{\tfrac{4}{3}}$}
	{\True $a^2$}
	\loigiai{
		Ta có $a^{\tfrac{5}{3}}\cdot a^{\tfrac{1}{3}}=a^{\tfrac{5}{3}+\tfrac{1}{3}}={a^2}$. 
	}
\end{ex}
\Closesolutionfile{ans}
\indapan{6}{ans/ans-0-B1}
\cauds
\Opensolutionfile{ans}[ans/ans-0-B1-DS]

\begin{ex}%[1D6N1-4]
	Với $a>0$, $b>0$, $m$, $n$ là các số thực tùy ý. Giả sử các biểu thức xuất hiện trong các công thức của mỗi mệnh đề đều có nghĩa.
	\choiceTF
	{\True $a^m\cdot a^n=a^{m+n}$}
	{\True$a^{-n}=\dfrac{1}{a^n}$ với $n$ nguyên dương}
	{\True Nếu $\left(\dfrac{5}{2}\right)^m>\left(\dfrac{5}{2}\right)^n$ thì $m>n$}
	{Nếu $\left(\dfrac{\pi}{4}\right)^m>\left(\dfrac{\pi}{4}\right)^{3m-2}$ thì $m<1$}
	\loigiai{
		\begin{itemchoice}
			\itemch \textbf{Đúng}\\
			Ta có $a^m\cdot a^n=a^{m+n}$ với $a>0$ và $m$, $n$ là các số thực tùy ý.
			\itemch \textbf{Đúng}\\
			Ta có $a^{-n}=\dfrac{1}{a^n}$ với $n$ nguyên dương.
			\itemch \textbf{Đúng}\\
			Vì $\dfrac{5}{2}>1$ nên nếu $\left(\dfrac{5}{2}\right)^m>\left(\dfrac{5}{2}\right)^n$ thì $m>n$.
			\itemch \textbf{Sai} \\
			Vì $0<\dfrac{\pi}{4}<1$ nên 
			\[\left(\dfrac{\pi}{4}\right)^m>\left(\dfrac{\pi}{4}\right)^{3m-2}\Leftrightarrow m<3m-2\Leftrightarrow m>1.\] 
		\end{itemchoice}
	}
\end{ex}

\begin{ex}%[1D6H1-1]
	Các câu sau đúng hay sai?
	\choiceTF
	{\True $\left[(0{,}2)^{\tfrac{3}{5}}\right]^{-5}=(0{,}2)^{-3}$}
	{\True $\left(5^{\tfrac{2}{3}}\right)^{3}=5^{2}$}
	{$\left(5^{\tfrac{2}{3}}\right)^{-3}+\left[(0{,}2)^{\tfrac{3}{5}}\right]^{-5}=5^m+5^n$ với $m$, $n$ là các số tự nhiên chẵn}
	{$\left(5^{\tfrac{2}{3}}\right)^{-3}+\left[(0{,}2)^{\tfrac{3}{5}}\right]^{-5}=K$ với $K$ chia hết cho 4}
	\loigiai{
		\begin{itemchoice}
			\itemch Đúng. $\left[(0{,}2)^{\tfrac{3}{5}}\right]^{-5}=(0{,}2)^{\tfrac{3}{5}\cdot(-5)}=(0{,}2)^{-3}$
			\itemch Đúng. $\left(5^{\tfrac{2}{3}}\right)^{3}=5^{{\tfrac{2}{3}}\cdot(3)}=5^{2}$.
			\itemch Sai. $\left(5^{\tfrac{2}{3}}\right)^{3}+\left[(0{,}2)^{\tfrac{3}{5}}\right]^{-5}=5^{2}+(0{,}2)^{-3}=5^{2}+\left(\dfrac{1}{5}\right)^{-3}=5^{2}+5^3\Rightarrow m=2$; $n=3$ là số lẻ. 
			\itemch Sai. $\left(5^{\tfrac{2}{3}}\right)^{3}+\left[(0{,}2)^{\tfrac{3}{5}}\right]^{-5}=5^{2}+(0{,}2)^{-3}=5^{2}+\left(\dfrac{1}{5}\right)^{-3}=5^{2}+5^3=150$.\\ 
			Mà $150$ không chia hết cho $4$. Nên $K$ không chia hết cho $4$.
	\end{itemchoice}}
\end{ex}


\begin{ex}%[1D6H1-1]
	Xét đúng sai các câu sau?
	\choiceTF
	{\True $81^{-0{,}75}+\left(\dfrac{1}{625}\right)^{-\tfrac{1}{4}}-\left(\dfrac{1}{32}\right)^{\tfrac{3}{5}}=3^m+5-2^n$, với $m+n=0$}
	{$\left(\dfrac{1}{625}\right)^{-\tfrac{1}{4}}=\left(5^{4}\right)^{-\tfrac{1}{4}}$}
	{$81^{-0{,}75}+\left(\dfrac{1}{625}\right)^{\tfrac{1}{4}}-\left(\dfrac{1}{32}\right)^{\tfrac{3}{5}}=-\dfrac{a}{b}\left(a, b \in \mathbb{N}^{*}\right)$ và $\dfrac{a}{b}$ là phân số tối giản, khi đó $a-b=52$}
	{\True $81^{-0{,}75}=\left(3^4\right)^{-\tfrac{3}{4}}$}
	\loigiai{
		\begin{itemchoice}
			\itemch Đúng. $81^{-0{,}75}+\left(\dfrac{1}{625}\right)^{-\tfrac{1}{4}}-\left(\dfrac{1}{32}\right)^{\tfrac{3}{5}}=\left(3^4\right)^{-\tfrac{3}{4}}+\left(5^{-4}\right)^{-\tfrac{1}{4}}-\left(2^{-5}\right)^{\tfrac{3}{5}}=3^{-3}+5-2^3\Rightarrow m=-3$; $n=3$.\\ $\Rightarrow m+n=0$.
			\itemch Sai. $\left(\dfrac{1}{625}\right)^{-\tfrac{1}{4}}=\left(5^{-4}\right)^{-\tfrac{1}{4}}$.
			\itemch Sai. $81^{-0{,}75}+\left(\dfrac{1}{625}\right)^{-\tfrac{1}{4}}-\left(\dfrac{1}{32}\right)^{\tfrac{3}{5}}=\left(3^4\right)^{-\tfrac{3}{4}}+\left(5^{-4}\right)^{-\tfrac{1}{4}}-\left(2^{-5}\right)^{\tfrac{3}{5}}=3^{-3}+5-2^3=-\dfrac{80}{27}\Rightarrow a =80$; $b=27$.\\
			$\Rightarrow a-b=52$.
			\itemch Đúng. $81^{-0{,}75}=\left(3^4\right)^{-\tfrac{3}{4}}$.
	\end{itemchoice}}
\end{ex}

\begin{ex}%[1D6N1-1]
	Cho biểu thức $9^{\tfrac{2}{5}} \cdot 27^{\tfrac{2}{5}}=A$ và $144^{\tfrac{3}{4}}: 9^{\tfrac{3}{4}}=B$, khi đó
	\choiceTF
	{\True $A-B=1$}
	{\True $144^{\tfrac{3}{4}}: 9^{\tfrac{3}{4}}=2^k$ thì $k=3$}
	{\True $9^{\tfrac{2}{5}} \cdot 27^{\tfrac{2}{5}}=(9 \cdot 27)^{\tfrac{2}{5}}$}
	{$9^{\tfrac{2}{5}} \cdot 27^{\tfrac{2}{5}}=3^k$ thì $k=3$}
	\loigiai{
		\begin{itemchoice}
			\itemch Đúng. Ta có $A-B=9-8=1$.
			\itemch Đúng. Ta có $144^{\tfrac{3}{4}}: 9^{\tfrac{3}{4}}=\left(\dfrac{144}{9}\right)^{\tfrac{3}{4}}=16^{\tfrac{3}{4}}=\left(2^4\right)^{\tfrac{3}{4}}=2^3=8=2^3\Rightarrow k=3$.
			\itemch Đúng. Ta có $9^{\tfrac{2}{5}} \cdot 27^{\tfrac{2}{5}}=(9 \cdot 27)^{\tfrac{2}{5}}$.
			\itemch Sai. Ta có $9^{\tfrac{2}{5}} \cdot 27^{\tfrac{2}{5}}=(9 \cdot 27)^{\tfrac{2}{5}}=\left(3^2 \cdot 3^3\right)^{\tfrac{2}{5}}=\left(3^5\right)^{\tfrac{2}{5}}=3^2=9=3^2\Rightarrow k =2.$
	\end{itemchoice}}
\end{ex}
\Closesolutionfile{ans}
\indapan{3}{ans/ans-0-B1-DS}

\Opensolutionfile{ans}[ans/ans-0-B1-KQ]
\caukq
\begin{ex}%[1D6H1-2]
	Cho $x$, $y$ là hai số nguyên thỏa mãn $3^x \cdot 6^y=\dfrac{2^{15} \cdot 6^{40}}{9^{50} \cdot 12^{25}}$. Tính $x y$.
	\shortans{$-450$}
	\loigiai{
		Ta có 	
		\[VT=3^x \cdot 6^y=3^x \cdot(2 \cdot 3)^y=2^y \cdot 3^{x+y}.\]
		Lại có \[VP=\dfrac{2^{15} \cdot 6^{40}}{9^{50} \cdot 12^{25}}=\dfrac{2^{15} \cdot(2 \cdot 3)^{40}}{\left(3^2\right)^{50} \cdot\left(2^2 \cdot 3\right)^{25}}=\dfrac{2^{15} \cdot 2^{40} \cdot 3^{40}}{3^{100} \cdot 2^{50} \cdot 3^{25}}=\dfrac{2^{55} \cdot 3^{40}}{2^{50} \cdot 3^{125}}=2^5 \cdot 3^{-85}.\]
		Suy ra \[2^y \cdot 3^{x+y}=2^5 \cdot 3^{-85} \Leftrightarrow \heva{&y=5 \\&x+y=-8 5} \Leftrightarrow \heva{&y=5 \\
			&x=-90.}\]
		Vậy $x \cdot y=5 \cdot(-90)=-450$.
		
	}
\end{ex}

\begin{ex}%[1D6V1-2]
	Cho biểu thức
	$$
	P=\left[\dfrac{\sqrt[3]{a^2 b}-\sqrt[3]{a b^2}}{\sqrt[3]{a^2}-2 \sqrt[3]{a b}+\sqrt[3]{b^2}}-\dfrac{a+b}{\sqrt[3]{a^2}-\sqrt[3]{b^2}}\right](\sqrt[6]{a}+\sqrt[6]{b})^{-1}+\sqrt[6]{a}, \text{ với }a>0, b>0, a \neq b.
	$$
	Tính $\dfrac{P^6}{b}$.
	\shortans{$1$}
	\loigiai{
		Ta có:
		$$
		\begin{aligned}
			& P=\left[\dfrac{\sqrt[3]{a^2 b}-\sqrt[3]{a b^2}}{\sqrt[3]{a^2}-2 \sqrt[3]{a b}+\sqrt[3]{b^2}}-\dfrac{a+b}{\sqrt[3]{a^2}-\sqrt[3]{b^2}}\right](\sqrt[6]{a}+\sqrt[6]{b})^{-1}+\sqrt[6]{a}\\
			& =\left[\dfrac{\sqrt[3]{a b}(\sqrt[3]{a}-\sqrt[3]{b})}{(\sqrt[3]{a}-\sqrt[3]{b})^2}-\dfrac{(\sqrt[3]{a}+\sqrt[3]{b})\left(\sqrt[3]{a^2}-\sqrt[3]{a b}+\sqrt[3]{b^2}\right)}{(\sqrt[3]{a}-\sqrt[3]{b})(\sqrt[3]{a}+\sqrt[3]{b})}\right](\sqrt[6]{a}+\sqrt[6]{b})^{-1}+\sqrt[6]{a}\\
			& =\left[\dfrac{\sqrt[3]{a b}}{\sqrt[3]{a}-\sqrt[3]{b}}-\dfrac{\sqrt[3]{a^2}-\sqrt[3]{a b}+\sqrt[3]{b^2}}{\sqrt[3]{a}-\sqrt[3]{b}}\right] \dfrac{1}{\sqrt[6]{a}+\sqrt[6]{b}}+\sqrt[6]{a}\\
			& =-\dfrac{(\sqrt[3]{a}-\sqrt[3]{b})^2}{\sqrt[3]{a}-\sqrt[3]{b}}\dfrac{1}{(\sqrt[6]{a}+\sqrt[6]{b})}+\sqrt[6]{a}\\
			& =\dfrac{\sqrt[3]{b}-\sqrt[3]{a}}{\sqrt[6]{a}+\sqrt[6]{b}}+\sqrt[6]{a}-\sqrt[6]{a}+\sqrt[6]{a}=\sqrt[6]{b}
		\end{aligned}
		$$
		Vậy $\dfrac{P^6}{b}=1$.
	}
\end{ex}

\begin{ex}%[1D6V1-3]
	Có bao nhiêu giá trị nguyên của tham số $m$ để hàm số $f(x)=\left(2x^2+mx+2\right)^{\tfrac{1}{2}}$ xác định với mọi $x\in \mathbb{R}$?\par
	\shortans{$7$}
	\loigiai{
		Do $\dfrac{1}{2}\notin \mathbb{Z}$ nên hàm số đã cho xác định $\Leftrightarrow 2x^2+mx+2>0$.\\
		Hàm số đã cho  xác định với mọi $x\in \mathbb{R}$ khi và chỉ khi
		$$2x^2+mx+2>0, \forall x \in \mathbb{R} \Leftrightarrow \Delta = m^2-16<0 \Leftrightarrow -4<m<4.$$
		Vì $m \in \mathbb{Z}$ nên $m\in \{-3;-2;-1;0;1;2;3\}$. Vậy có $7$ giá trị nguyên của $m$ thỏa yêu cầu bài toán.
	}
\end{ex}

\begin{ex}%[1D6V1-1]
	Anh Toàn được tuyển dụng vào một công ty đầu năm 2013. Công ty trả lương cho anh theo nguyên tắc: Lương khởi điểm anh nhận là $6$ triệu đồng/ tháng và cứ sau $3$ năm công ty lại tăng lương cho anh thêm $25\%$ số lương đang hưởng. Hiện nay (năm 2024) anh đang được hưởng lương là bao nhiêu triệu đồng một tháng? (Kết quả làm tròn đến hàng phần mười)
	\par
	\shortans[0]{$11{,}7$}
	\loigiai{
		Tính từ năm 2013 đến năm 2024, anh Toàn đã được $3$ lần tăng lương.\\
		Lương của anh Toàn sau lần tăng đầu tiên là $L_1=6\cdot 1{,}25$ triệu.\\
		Lương của anh Toàn sau lần tăng thứ hai là $L_2=L_1+25\%\cdot L_1=L_1\cdot 1{,}25=6\cdot (1{,}25)^2$ triệu.\\
		Lương của anh Toàn sau lần tăng thứ ba là $L_3=L_2+25\cdot \%L_2=6\cdot (1{,}25)^3\approx 11{,}7 $ triệu.\\
		Vậy lương của anh Toàn hiện đang hưởng là $11{,}7$ triệu mỗi tháng.
	}
\end{ex}

\begin{ex}%[1D6V1-1]
	Cho biểu thức $A=(\sqrt{2}-1)^{2 x}+(3+2 \sqrt{2})^x$. Tính $A$ khi $(\sqrt{2}+1)^x=2$.
	\shortans[0]{$4{,}25$}
	\loigiai{
		Ta có  $\sqrt{2}-1=\dfrac{(\sqrt{2}-1)(\sqrt{2}+1)}{\sqrt{2}+1}=\dfrac{1}{\sqrt{2}+1}=(\sqrt{2}+1)^{-1}$.\\
		Mà $3+2 \sqrt{2}=(\sqrt{2})^2+2 \sqrt{2}+1=(\sqrt{2}+1)^2$.\\
		Do đó 
		\begin{eqnarray*}
			A&=&\left[(\sqrt{2}+1)^{-1}\right]^{2 x}+\left[(\sqrt{2}+1)^2\right]^x\\
			&=&(\sqrt{2}+1)^{-2 x}+(\sqrt{2}+1)^{2 x}\\
			&=&\left[(\sqrt{2}+1)^x\right]^{-2}+\left[(\sqrt{2}+1)^x\right]^2.
		\end{eqnarray*}
		Vậy $A = 4{,}25$.
	}
\end{ex}
\begin{ex}%[1D6H1-1]
	Cho biểu thức $A=3^{2 x-1} \cdot\left(\dfrac{1}{3}\right)^{2 x-1}+9^{x+1}$.
	Tính $A$ khi $3^x=2$.
	\shortans[0]{$37$}
	\loigiai{
		Ta có $A=\left(3 \cdot \dfrac{1}{3}\right)^{2 x-1}+9 \cdot 9^x=1+9\left(3^x\right)^2 = 1+9\left(2\right)^2=37$.
	}
\end{ex}
\Closesolutionfile{ans}
\indapan{6}{ans/ans-0-B1-KQ}